% !TEX root = ../main.tex
% 光子晶体面发射激光器偏振调控技术汇总
\begin{figure}[htbp]
  \centering
  \begin{tikzpicture}[x=0.95cm,y=0.95cm,>=Latex,font=\small]
    \node[font=\Large\bfseries] at (0,5.3) {光子晶体面发射激光器偏振调控技术汇总};

    % (a) 椭圆孔
    \begin{scope}[shift={(-4.6,2.5)}]
      \draw[thick] (-2.3,-1.3) rectangle (2.3,1.3);
      \foreach \x in {-1.7,-0.8,0.1,1.0,1.9} {
        \foreach \y in {-0.9,-0.3,0.3,0.9} {
          \draw[fill=gray!25] (\x,\y) ellipse (0.15 and 0.24);
        }
      }
      \draw[->,very thick,red] (-2.1,-1.7) -- (2.1,-1.7);
      \node[red,below] at (0,-2.25) {(a) 椭圆孔各向异性};
      \node[above] at (0,1.45) {破坏 $C_4$ 对称性};
    \end{scope}

    % (b) 矩形晶格
    \begin{scope}[shift={(4.6,2.5)}]
      \draw[thick] (-2.3,-1.3) rectangle (2.3,1.3);
      \foreach \x in {-1.6,-0.7,0.2,1.1,2.0} {
        \foreach \y in {-0.8,0.0,0.8} {
          \draw[fill=gray!25] (\x,\y) circle (0.14);
        }
      }
      \draw[<->] (-2.45,-1.0) -- (-2.45,1.0) node[midway,left] {$a_y$};
      \draw[<->] (-1.8,-1.55) -- (1.8,-1.55) node[midway,below] {$a_x$};
      \node[above] at (0,1.45) {$a_x\neq a_y$};
      \node[blue,below] at (0,-2.25) {(b) 矩形晶格};
    \end{scope}

    % (c) 狭缝型
    \begin{scope}[shift={(-4.6,-1.4)}]
      \draw[thick] (-2.3,-1.3) rectangle (2.3,1.3);
      \fill[gray!25] (-2.3,-1.3) rectangle (2.3,1.3);
      \fill[white] (-0.35,-1.3) rectangle (0.35,1.3);
      \draw[thick] (-0.35,-1.3) rectangle (0.35,1.3);
      \foreach \y in {-0.85,-0.25,0.35,0.95} {
        \draw[fill=white] (-1.35,\y) circle (0.18);
        \draw[fill=white] (1.35,\y) circle (0.18);
      }
      \draw[<->,very thick,red] (-0.28,0.1) -- (0.28,0.1);
      \node[red,below] at (0,-2.15) {(c) 狭缝型光子晶体};
      \node[above] at (0,1.45) {狭缝内场增强};
    \end{scope}

    % (d) 集成偏振栅
    \begin{scope}[shift={(4.6,-1.4)}]
      \draw[thick] (-2.3,-1.3) rectangle (2.3,1.3);
      \foreach \x in {-1.9,-1.5,...,1.7} {
        \fill[gray!65] (\x,-1.3) rectangle ++(0.18,2.6);
      }
      \draw[->,very thick,green!60!black] (-3.1,0.12) -- (-2.28,0.12);
      \node[left,green!60!black,align=right] at (-3.1,0.12) {\scriptsize 横电入射};
      \draw[->,thick,red,dashed] (-3.1,0.62) -- (-2.28,0.62);
      \node[left,red,align=right] at (-3.1,0.62) {\scriptsize 横磁入射};
      \draw[very thick,red] (2.28,0.55) -- (2.85,0.55);
      \node[red,right] at (2.85,0.55) {被吸收};
      \node[blue,below] at (0,-2.15) {(d) 集成线栅偏振片};
      \node[above] at (0,1.45) {选择性偏振损耗};
    \end{scope}

    \node[draw,align=center,text width=11.2cm] at (0,-4.95)
      {\textbf{典型偏振消光比：}
      椭圆孔：约 $8$--$12$ 分贝，矩形晶格：约 $10$--$15$ 分贝，狭缝型：高于 $15$ 分贝，线栅偏振片：高于 $20$ 分贝。\\
      核心思路：引入可控各向异性，使目标偏振的阈值低于正交偏振。};
  \end{tikzpicture}
  \caption{光子晶体面发射激光器中四类常用偏振调控方法对比。}
  \label{fig:ch08_polarization_control}
\end{figure}
